\section{Simulaciones 1 esfera}

\subsection{Régimen de Rayleigh: $x\ll1$}
\subsubsection{Barrido de longitudes de onda en el rango visible}
En esta sección se simula una esfera de oro aislada de radio $a = 20\,\text{nm}$
en el rango visible del espectro electromagnético, desde $\lambda = 400\,\text{nm}$
hasta $\lambda = 800\,\text{nm}$ con un paso de $5\,\text{nm}$.
Para cada longitud de onda, el índice de refracción complejo del oro se obtiene
mediante interpolación de los datos experimentales de Johnson \& Christy
\cite{johnson1972optical}.

La figura~\ref{fig:qext_qabs_qsca_r20sw} muestra las eficiencias de extinción, absorción y
dispersión en función de la longitud de onda.

\begin{figure}[H]
    \centering
    \includegraphics[width=0.9\textwidth]{../../Simulaciones/Definitivo/simulacion-de-referencia/rayleigh/r20sw/plots/Efficiencies_-_Unpol.png}
    \caption{Eficiencias de extinción, absorción y dispersión de una esfera de oro de radio $20 \,n$m frente a la longitud de onda.}
    \label{fig:qext_qabs_qsca_r20sw}
\end{figure}

Las tres cantidades presentan un máximo en torno a 
$\lambda \approx 500 - 520 \,\text{nm}$,
a partir del cual decrecen al aumentar la longitud de onda.
En todo el rango estudiado, la eficiencia de absorción $Q_\text{abs}$ domina
claramente sobre la de dispersión $Q_\text{sca}$, con valores de $Q_\text{sca}$
aproximadamente treinta veces menores que $Q_\text{abs}$.
Esto indica que la mayor parte de la energía extraída del haz incidente es
absorbida por la partícula y convertida en calor, mientras que la contribución
de la dispersión es comparativamente despreciable. 

La figura~\ref{fig:S11_small} muestra el elemento $S_{11}(\theta)$ de la
matriz de dispersión en función del ángulo polar $\theta$, para varias
longitudes de onda del rango visible.

\begin{figure}[H]
    \centering
    \includegraphics[width=0.7\textwidth]{../../Simulaciones/Definitivo/simulacion-de-referencia/rayleigh/r20sw/plots/S11.png}
    \caption{Elemento $S_{11}(\theta)$ de la matriz de dispersión para una esfera de oro de radio $20\,\text{nm}$ frente al ángulo polar $\theta$ en el campo lejano.}
    \label{fig:S11_small}
\end{figure}



El patrón angular es prácticamente simétrico respecto a $\theta = 90°$,
sin preferencia neta por la dispersión hacia delante ni hacia atrás.
En consecuencia, el parámetro de asimetría $g = \langle\cos\theta\rangle$
toma un valor próximo a cero.
Asimismo, las curvas correspondientes a las distintas longitudes de onda
se superponen casi exactamente, lo que indica que el patrón angular es
prácticamente independiente de $\lambda$ en este régimen.


\begin{figure}[htbp]
    \centering
    \includegraphics[width=0.7\textwidth]{../../Simulaciones/Definitivo/simulacion-de-referencia/rayleigh/r20sw/plots/ALL.png}
    \caption{Grado de polarización total de la luz dispersada por una esfera de oro de radio $20\,\text{nm}$ frente al ángulo polar $\theta$ en el campo lejano.}
    \label{fig:ALL_small}
\end{figure}



La figura~\ref{fig:ALL_small} muestra el grado de polarización lineal
$\text{DoLP}(\theta)$ y el grado de polarización circular $\text{DoCP}(\theta)$
superpuestos en función del ángulo de dispersión, para varias 
longitudes de onda del rango visible.


El $\text{DoCP}$ es nulo en todo el rango angular y para todas las longitudes
de onda consideradas, lo que indica que la luz dispersada no adquiere
polarización circular. Como consecuencia, el grado de polarización total
$\text{DoP} = \sqrt{\text{DoLP}^2 + \text{DoCP}^2}$ coincide exactamente
con el $\text{DoLP}$, y ambas curvas se superponen en la figura.
El $\text{DoLP}$ presenta un máximo en $\theta = 90°$, donde la luz
dispersada está completamente polarizada linealmente, y se anula en las
direcciones axiales $\theta = 0°$ y $\theta = 180°$.
Al igual que el patrón angular de $S_{11}$, el $\text{DoLP}$ es
prácticamente independiente de la longitud de onda en este régimen.




\begin{figure}[htbp]
    \centering

    \begin{subfigure}[b]{0.31\textwidth}
        \centering
        \includegraphics[width=\textwidth]{../../Simulaciones/Definitivo/simulacion-de-referencia/rayleigh/r20sw/plots/nf500.png}
        \caption{$\lambda = 500\,\text{nm}$.}
        \label{fig:nearfield_small_500}
    \end{subfigure}
    \hfill
    \begin{subfigure}[b]{0.31\textwidth}
        \centering
        \includegraphics[width=\textwidth]{../../Simulaciones/Definitivo/simulacion-de-referencia/rayleigh/r20sw/plots/nf520.png}
        \caption{$\lambda = 520\,\text{nm}$.}
        \label{fig:nearfield_small_520}
    \end{subfigure}
    \hfill
    \begin{subfigure}[b]{0.31\textwidth}
        \centering
        \includegraphics[width=\textwidth]{../../Simulaciones/Definitivo/simulacion-de-referencia/rayleigh/r20sw/plots/nf540.png}
        \caption{$\lambda = 540\,\text{nm}$.}
        \label{fig:nearfield_small_540}
    \end{subfigure}

    \caption{Distribución espacial de la intensidad del campo eléctrico en el campo cercano 
            para tres longitudes de onda.}
    \label{fig:nearfield_small}
\end{figure}


La figura~\ref{fig:nearfield_small} muestra la distribución espacial de la
intensidad del campo eléctrico en las proximidades de la esfera para tres
longitudes de onda: $\lambda = 500\,\text{nm}$, $\lambda = 520\,\text{nm}$
y $\lambda = 540\,\text{nm}$, correspondientes al entorno del máximo de
$Q_\text{sca}$ observado en la figura~\ref{fig:qext_qabs_qsca_r20sw}.
El campo cercano se evalúa en un plano que contiene el eje de propagación
de la onda incidente, y se representa normalizado respecto a la amplitud
del campo incidente.

Al tratarse de una esfera aislada, no existen efectos de acoplamiento entre
partículas y el campo en la región exterior refleja únicamente la
perturbación introducida por la presencia de la esfera sobre la onda plana
incidente.

